\chapter{JSON API, CORS, and browser client integration}

RWLang does not force you to choose between classic server-rendered HTML and a JSON API. The same application can contain pages returning \code{Html}, actions processing forms, and typed JSON endpoints. This is useful both when a small JavaScript component needs some data and when a separate frontend application uses the RWLang server as an API.

This chapter cleanly separates three operating models:

\begin{enumerate}
  \item \textbf{same-origin HTML + API}: the browser loads the page and calls the API from the same origin;
  \item \textbf{cross-origin public API}: another origin calls the API, but there is no session cookie;
  \item \textbf{cross-origin session auth}: a separate frontend origin uses cookie-based sessions, so CORS and CSRF policy are both required.
\end{enumerate}

The most important rule is: \textbf{CORS is neither authentication nor CSRF protection}. CORS is a browser-side access policy; authentication decides who the user is, while CSRF protection prevents an authenticated session from being abused for unauthorized state changes from another origin.

\section{The first typed JSON endpoint}

A JSON endpoint handler returns \code{Json}. \code{json(...)} converts a supported runtime value into a JSON representation.

\begin{lstlisting}[language=RWLang]
page fn statusApi(ctx: PageContext) -> Result<Json, PageError> {
    let healthy = true;
    return Ok(json(healthy));
}

route statusApi GET "/api/status" => statusApi;
\end{lstlisting}

A successful response is JSON, for example:

\Needspace{9\baselineskip}
\begin{lstlisting}
HTTP/1.1 200 OK
Content-Type: application/json; charset=utf-8

true
\end{lstlisting}

\code{json(value)} is not an arbitrary object serializer. Supported values are constrained: scalar values, model records, optional models, and lists of models. This matches RWLang's typed-boundary philosophy: the response shape does not emerge from an arbitrary runtime-changing data structure.

\section{JSON lists from the database}

Typed SQL and JSON responses connect directly. The query still returns an exact model shape, and the JSON layer represents that typed value.

\begin{lstlisting}[language=RWLang]
model Product {
    id: Int
    name: String
    price: Int
}

query fn listProducts(db: Db, limit: Int, offset: Int)
    -> Result<List<Product>, DbError> sql {
    SELECT id, name, price
    FROM products
    ORDER BY id
    LIMIT :limit OFFSET :offset
}

page fn productsApi(ctx: PageContext, db: Db,
                    limit: Int, offset: Int)
    -> Result<Json, PageError> {
    let products = listProducts(db, limit, offset)?;
    return Ok(json(products));
}

route productsApi GET "/api/products"
    query limit<Int> offset<Int>
    validate limit range 1 100 offset range 0 1000000
    => productsApi;
\end{lstlisting}

The important point is not the JSON but the fact that input remains typed and validated just as it is for an HTML route. A client cannot bypass route policy by supplying an unbounded \code{limit}.

\Needspace{17\baselineskip}
\section{Typed JSON POST body}

A route declares \code{json} input mode for a JSON request body.

\begin{lstlisting}[language=RWLang]
action fn echoApi(ctx: ActionContext,
                  name: String,
                  age: Int,
                  active: Bool)
    -> Result<Json, PageError> {
    return Ok(json(name));
}

route echoApi POST "/api/echo"
    json name<String> age<Int> active<Bool>
    validate name length 1 100 age range 0 150
    auth user
    => echoApi;
\end{lstlisting}

A valid request is:

\begin{lstlisting}
POST /api/echo HTTP/1.1
Content-Type: application/json
Accept: application/json
X-CSRF-Token: <session CSRF token>

{"name":"Alice","age":42,"active":true}
\end{lstlisting}

The JSON input schema is \textbf{closed}. Unknown or missing keys, duplicate keys, wrong scalar types, and unsupported nested objects, arrays, floats, or \code{null} values are errors. These are request-contract violations and therefore return \code{400 Bad Request}. A missing or incorrect \code{Content-Type} on a JSON route returns \code{415 Unsupported Media Type}. Field and request-size limits apply just as they do to other request bodies.

The \code{json}, \code{form}, and \code{upload} body modes are mutually exclusive on the same POST route, and the compiler enforces this. JSON expects \code{application/json}. Simple forms use URL-encoded form data; uploads use multipart form data. Each endpoint should have one unambiguous request contract.

\section{Content negotiation and 406}

The client can use the \code{Accept} header to state what representation it can receive. For a JSON route, supported examples include:

\begin{lstlisting}
Accept: application/json
Accept: application/*
Accept: */*
\end{lstlisting}

If a handler would return JSON but the client accepts only \code{text/html}, the server returns \code{406 Not Acceptable}. This is preferable to silently returning a representation the client explicitly rejected.

\section{API errors}

On JSON routes, runtime-classified application errors use a simple common envelope: a stable \code{error} code, while detailed diagnosis remains in server-side logs. Clients should not parse the English text of error messages; they should use the HTTP status and stable error code.

Once a route is known to return JSON, errors use a JSON envelope, for example:

\begin{lstlisting}
HTTP/1.1 401 Unauthorized
Content-Type: application/json; charset=utf-8

{"error":"unauthorized"}
\end{lstlisting}

Typical error codes include:

\begin{itemize}
  \item \code{bad_request};
  \item \code{unauthorized};
  \item \code{forbidden};
  \item \code{csrf_failed};
  \item \code{unsupported_media_type};
  \item \code{database_unavailable};
  \item \code{resource_limit};
  \item \code{internal_error}.
\end{itemize}

This is especially important for browser clients. An authenticated JSON route returns \code{401} to an anonymous request rather than an HTML login redirect. The frontend can therefore decide based on status code without having to interpret HTML as an API error.

\section{Same-origin: the default and simplest model}

If HTML and API are served from the same scheme + host + port combination, they are same-origin from the browser's point of view.

\begin{lstlisting}
https://app.example/
https://app.example/api/products
\end{lstlisting}

No CORS configuration is normally needed. A simple client is:

\begin{lstlisting}
const response = await fetch("/api/products?limit=20&offset=0", {
  headers: { Accept: "application/json" }
});

if (!response.ok) {
  throw new Error(`HTTP ${response.status}`);
}

const products = await response.json();
\end{lstlisting}

Prefer this architecture when there is no real business need for a separate frontend origin. It means fewer trust boundaries, a simpler cookie policy, and less deployment configuration.

\section{CORS: exact-origin allowlist}

Cross-origin browser access is denied by default. Trusted server configuration must enumerate exact origins; there is no wildcard origin.

\begin{lstlisting}
[web]
cors_origins = ["https://frontend.example"]
cors_allow_credentials = false
\end{lstlisting}

The CLI equivalent is:

\begin{lstlisting}
rwlang-server --cors-origin https://frontend.example ...
\end{lstlisting}

Use multiple explicit allowlist entries for multiple origins. Do not build a policy that blindly reflects the request's \code{Origin} value.

\section{Preflight}

Before some cross-origin requests, the browser sends an \code{OPTIONS} preflight. RWLang allows the preflight only when:

\begin{itemize}
  \item the \code{Origin} is configured;
  \item the target GET/POST route exists;
  \item the requested request headers are supported.
\end{itemize}

Supported browser-API headers include \code{Content-Type}, \code{Accept}, and \code{X-CSRF-Token}. Preflight is therefore not a generic ``allow everything'' OPTIONS endpoint; it is part of the real route and origin policy.

\section{Credentialed CORS with a session cookie}

When a frontend on a different origin uses RWLang session authentication, credentialed CORS is required:

\begin{lstlisting}
[web]
cors_origins = ["https://frontend.example"]
cors_allow_credentials = true
\end{lstlisting}

In production this is allowed only over HTTPS. With this configuration the session cookie uses \code{SameSite=None; Secure}, because the browser must send it on cross-site requests.

The JavaScript fetch call also needs explicit credential mode:

\begin{lstlisting}
const response = await fetch("https://api.example/api/account", {
  credentials: "include",
  headers: { Accept: "application/json" }
});
\end{lstlisting}

\code{credentials: "include"} does not grant permission. It only tells the browser that matching cookies may be sent for a request allowed by the CORS policy.

\section{CSRF on a JSON API}

With cookie-based sessions, the browser manages the session cookie automatically, so state-changing JSON requests need CSRF protection just like HTML forms.

For form POST, the token is a body field:

\begin{lstlisting}
_csrf=<token>
\end{lstlisting}

For JSON POST, use a header:

\begin{lstlisting}
X-CSRF-Token: <token>
\end{lstlisting}

For a separate frontend origin, an authenticated token endpoint can be defined:

\begin{lstlisting}[language=RWLang]
page fn csrfApi(ctx: PageContext) -> Result<Json, PageError> {
    return Ok(json(csrfToken));
}

route csrfApi GET "/api/csrf" auth user => csrfApi;
\end{lstlisting}

The token belongs to the same session. A CORS allowlist does not replace this check.

\section{Complete browser POST flow}

A session-auth frontend running from a separate \code{frontend.example} origin typically does the following:

\begin{enumerate}
  \item the user authenticates and the browser receives the session cookie;
  \item the frontend fetches the CSRF token belonging to that session;
  \item POST is sent with \code{credentials: "include"};
  \item the JSON body conforms to the route's closed typed schema;
  \item \code{X-CSRF-Token} contains the session token;
  \item the runtime verifies origin, authentication, CSRF, and input validation;
  \item only then does the handler perform the business operation.
\end{enumerate}

Example client:

\begin{lstlisting}
const csrf = await fetch("https://api.example/api/csrf", {
  credentials: "include",
  headers: { Accept: "application/json" }
}).then(async r => {
  if (!r.ok) throw new Error(`CSRF HTTP ${r.status}`);
  return r.json();
});

const result = await fetch("https://api.example/api/echo", {
  method: "POST",
  credentials: "include",
  headers: {
    "Content-Type": "application/json",
    Accept: "application/json",
    "X-CSRF-Token": csrf
  },
  body: JSON.stringify({ name: "Alice", age: 42, active: true })
}).then(async r => {
  if (!r.ok) throw new Error(`API HTTP ${r.status}`);
  return r.json();
});
\end{lstlisting}

\section{Reverse proxy and origin}

Behind Apache or Nginx, public scheme/host and trusted-proxy configuration must agree. The application must not decide what origin is safe based on arbitrary forwarded headers supplied by the client. The installation chapter explains the trusted-proxy boundary; the full \code{server.toml} key reference appears near the end of the book.

Always place the \textbf{frontend origin as seen by the browser} in the CORS allowlist, for example:

\begin{lstlisting}
https://frontend.example
\end{lstlisting}

not an internal proxy or container address.

\section{What not to do}

\begin{itemize}
  \item Do not enable CORS merely because ``this is an API''.
  \item Do not reflect arbitrary request \code{Origin} values.
  \item Do not treat CORS as authentication or CSRF protection.
  \item Do not send session credentials over HTTP in production.
  \item Do not mix form, JSON, and upload body modes on the same POST route.
  \item Do not leave list-query parameters unbounded; API input is validated input too.
  \item Do not build JSON-client logic around an HTML login redirect; handle 401/403 statuses.
\end{itemize}

\section{Which architecture should you choose?}

\begin{center}
\begin{tabularx}{\textwidth}{@{}L{0.27\textwidth}Y@{}}
\toprule
\textbf{Situation} & \textbf{Recommended model} \\
\midrule
server-rendered page + small JS & same-origin HTML + JSON API \\
SPA on the same domain & same-origin API, no CORS \\
separate public frontend, no session & exact-origin CORS without credentials \\
separate frontend + session auth & exact-origin credentialed CORS + HTTPS + CSRF \\
server-to-server client & CORS is irrelevant; auth and network policy matter \\
\bottomrule
\end{tabularx}
\end{center}

\section{Runnable example}

The complete repository example is \path{examples/json-api/app.rw}.

Try it first from the same origin. Introduce a separate frontend origin and CORS only when the deployment genuinely requires it. This keeps it clear which problem is solved by the JSON API and which is solved by the browser's cross-origin policy.
